\section{A Posteriori Evaluation}
\label{sec:aposteriori}

To assess whether improved offline accuracy translates to better coupled CFD predictions, we evaluate all models within the solver on three canonical flow configurations.
Each model uses the same spatial discretization, time integrator, and Poisson solver; only the turbulence closure differs.

\subsection{Channel flow at $\Reytau = 180$}

\subsubsection{Setup}

\todo{Specify grid ($128 \times 96 \times 128$), domain size ($4\pi \times 2 \times 2\pi$), $\Rey$ (or $\nu$, $dp/dx$), time integrator (RK3), number of steps, CFL, stretch parameter. Reference: MKM DNS \citep{MKM1999}.}

\subsubsection{Mean velocity profiles}

\todo{Plot $u^+(y^+)$ for all models vs.\ MKM DNS data. Report the $L_2$ error in $u^+$ for each model. Include the law of the wall ($u^+ = y^+$, log law $u^+ = (1/\kappa)\ln y^+ + B$) for reference.}

\subsubsection{Reynolds stress profiles}

\todo{Plot $-\langle u'v' \rangle / u_\tau^2$ vs.\ $y/\delta$ for all models vs.\ MKM DNS. Report component-wise $L_2$ errors. This is the most discriminating metric for tensor-predicting models (TBNN, TBRF) vs.\ scalar models (MLP, SST).}

\subsubsection{Predicted $\Reytau$}

\todo{Report the friction Reynolds number $\Reytau = u_\tau \delta / \nu$ achieved by each model. MKM target: $\Reytau = 178.12$.}

\subsection{Periodic hills at $\Rey_H = 5600$}

\subsubsection{Setup}

\todo{Specify grid, domain, $\Rey$, hill geometry (IBM), reference data \citep{Breuer2009}. Note that periodic hills at different slope angles were in the training set---this tests generalization to a different Reynolds number.}

\subsubsection{Velocity profiles}

\todo{Plot streamwise velocity $u/U_b$ vs.\ $y/H$ at key $x$-stations ($x/H = 0.5, 2, 4, 6, 8$) for all models vs.\ Breuer LES. These stations capture the recirculation zone, reattachment, and recovery regions.}

\subsubsection{Separation and reattachment}

\todo{Report separation and reattachment $x$-locations for each model. Breuer reference: separation at $x/H \approx 0.2$, reattachment at $x/H \approx 4.7$ for $\Rey_H = 5600$.}

\subsubsection{Separation bubble visualization}

\todo{Contour plots or streamlines of the mean flow field showing the separation bubble for selected models (baseline, SST, TBNN) vs.\ reference.}

\subsection{Cylinder at $\Rey = 100$}

\subsubsection{Setup}

\todo{Specify grid ($128 \times 64 \times 64$), domain, cylinder diameter and location (IBM), reference values. Note that the cylinder geometry was not in the McConkey training set---this tests generalization to an entirely unseen geometry.}

\subsubsection{Drag and lift coefficients}

\todo{Report time-averaged drag coefficient $C_d$, RMS lift coefficient $C_{l,\mathrm{rms}}$, and Strouhal number $\mathrm{St}$ for each model. Established reference values: $C_d \approx 1.35$, $C_{l,\mathrm{rms}} \approx 0.34$, $\mathrm{St} \approx 0.165$.}

\subsubsection{Wake profiles}

\todo{Plot streamwise velocity and Reynolds stress profiles in the near wake ($x/D = 1, 2, 5$) for selected models vs.\ reference data.}

\subsection{Grid sensitivity}

\todo{Run the TBNN and SST models on 4 grids of increasing resolution for the channel flow case to demonstrate grid convergence. Report the $L_2$ error in $u^+(y^+)$ vs.\ grid spacing.}
